\chapter{System Requirements and Problem Analysis}

The Mini Payment Gateway is positioned as an implemented MVP rather than a full
payment platform. Its purpose is to support merchant onboarding, QR-based
payment creation, asynchronous provider result processing, full refunds,
webhook delivery, operational investigation, and merchant self-visibility
through a separate dashboard experience.

\section{Problem Context}

The gateway solves two connected problems. First, merchants need a small but
reliable payment interface for creating QR payments, receiving asynchronous
provider outcomes, querying transaction status, and requesting full refunds.
Second, internal operators need visibility and recovery tools because payment
systems are asynchronous by nature: callbacks can arrive late, webhook delivery
can fail, and merchant readiness must be enforced before money-flow APIs are
usable.

The dashboard work revealed an architectural constraint that strongly shaped the
project. A merchant-facing dashboard cannot safely be implemented as frontend
screens alone. It requires a separate portal authentication model, merchant-
scoped read APIs, a dedicated merchant user table, and internal provisioning
flows. Without these boundaries, the system would either misuse merchant API
credentials for human login or risk cross-merchant data exposure.

\section{Implemented Scope And Non-Goals}

\begin{table}[H]
\centering
\caption{Implemented system scope}
\begin{tabularx}{\textwidth}{L{0.24\textwidth}Y}
\toprule
\textbf{Area} & \textbf{Implemented capability} \\
\midrule
Merchant lifecycle & Create merchants, submit and review onboarding evidence, create or rotate credentials, activate, suspend, and disable merchants. \\
Payments & Create dynamic QR payments, enforce semantic idempotency, query by transaction id or order id, and process provider payment result callbacks. \\
Refunds & Create full refunds, enforce refund rules, query refund state, and process provider refund result callbacks. \\
Webhooks & Create durable signed outbound webhook events, store delivery attempts, apply retry policy, and support internal manual retry. \\
Ops Dashboard & Internal operator interface for auth, RBAC, merchant workflows, evidence inspection, reconciliation, audit, internal users, and portal-user provisioning. \\
Merchant Dashboard & Merchant-scoped read-only interface for overview, analytics, payments, refunds, webhooks, profile, credentials, and password change. \\
Testing & Backend unit, route, and E2E tests; frontend typecheck, test, and build; compose validation; and browser smoke procedures. \\
\bottomrule
\end{tabularx}
\end{table}

The following items are intentionally out of scope for the implemented MVP:

\begin{itemize}
  \item settlement ledger and accounting workflows;
  \item disputes or chargebacks;
  \item partial refunds;
  \item multi-provider routing;
  \item merchant self-service onboarding;
  \item merchant profile editing from the Merchant Dashboard;
  \item production-grade notification, TLS, and observability hardening.
\end{itemize}

\section{Functional Requirements}

\begin{longtable}{L{0.24\textwidth}L{0.30\textwidth}L{0.36\textwidth}}
\caption{Functional requirements}\\
\toprule
\textbf{Requirement group} & \textbf{Rule} & \textbf{Implemented decision} \\
\midrule
\endfirsthead
\toprule
\textbf{Requirement group} & \textbf{Rule} & \textbf{Implemented decision} \\
\midrule
\endhead
Merchant lifecycle & Merchant statuses are \texttt{PENDING\_REVIEW}, \texttt{ACTIVE}, \texttt{REJECTED}, \texttt{SUSPENDED}, and \texttt{DISABLED}. & Onboarding approval is stored on \texttt{merchant\_onboarding\_cases}; operational status remains on \texttt{merchants}. \\
Payment creation & One active payment is allowed per merchant and order. & A \texttt{PENDING} payment is reused only when the duplicate request is semantically identical; success blocks a new attempt. \\
Payment states & Persisted states are \texttt{PENDING}, \texttt{SUCCESS}, \texttt{FAILED}, and \texttt{EXPIRED}. & Provider callbacks finalize pending payments, while late or mismatched evidence is routed to reconciliation. \\
Refunds & Only full refunds are implemented; refund window is seven days from \texttt{paid\_at}. & Refund records use \texttt{REFUND\_PENDING}, \texttt{REFUNDED}, and \texttt{REFUND\_FAILED}. Partial refunds are rejected by design. \\
Webhooks & HTTP 2xx means delivered; retry schedule is 1, 5, and 15 minutes after the first failure. & Delivery attempts are persisted; exhausted events become \texttt{FAILED}; internal Ops can manually retry eligible failed events. \\
Dashboards & Ops and Merchant users require different surfaces. & Ops Dashboard is internal and mutable; Merchant Dashboard is read-only except password change. \\
\bottomrule
\end{longtable}

\section{Technical Requirements}

The backend protects each boundary with the correct authentication mechanism.
Merchant backend integrations use HMAC headers:
\texttt{X-Merchant-Id}, \texttt{X-Access-Key}, \texttt{X-Timestamp}, and
\texttt{X-Signature}. Internal users authenticate through an HttpOnly session
cookie. Merchant portal users authenticate through a separate HttpOnly session
cookie with merchant scope derived from the authenticated portal account.

The system also enforces persistence rules where practical:

\begin{itemize}
  \item one active credential per merchant;
  \item one pending payment per merchant and order;
  \item one logical refund per merchant and refund id;
  \item at most one successful full refund per payment;
  \item positive monetary values and non-negative attempt counters;
  \item merchant-scoped analytics indexes for payment, refund, and webhook
  records.
\end{itemize}

\section{Operational Requirements}

The gateway must remain supportable by internal operators. Every important
manual action carries an audit reason, and audit rows capture the event code,
entity identity, actor type, actor id, before and after snapshots, and
timestamp. Sensitive fields such as raw credential secrets are masked in
responses and audit snapshots.

The merchant-facing dashboard must help merchants inspect their own data without
seeing internal reconciliation workflows, raw credential secrets, or data owned
by other merchants. This requirement directly drives the dedicated merchant
portal authentication model and the rule that Merchant Dashboard APIs never
accept a client-supplied \texttt{merchant\_id} for read scope after login.
